\section*{Chapter \ref{ch:s1:residual-and-principal-component-stat-arb} --- Residual and Principal-Component Stat Arb}

\begin{solution}{exo:s1:residual-and-principal-component-stat-arb:1}
$\kappa = -252 \ln 0.9 = 26.6$ a year; the mean-reversion time is $252/\kappa = 9.5$ days.
\end{solution}

\begin{solution}{exo:s1:residual-and-principal-component-stat-arb:2}
$m = 0.002/0.1 = 0.02$; $\sigma_{\mathrm{eq}} = \sqrt{10^{-4}/(1 - 0.81)} = 2.29\%$; $s = -0.02/0.0229 = -0.87$. No: $|s| < 1.25$.
\end{solution}

\begin{solution}{exo:s1:residual-and-principal-component-stat-arb:3}
$\kappa = -252 \ln b > 252/30$ is $\ln b < -1/30$, that is $b < e^{-1/30} = 0.9672$.
\end{solution}

\begin{solution}{exo:s1:residual-and-principal-component-stat-arb:4}
The weight traded is twice the turnover: $2 \times 0.37 \times 0.0005 \times 252 = 9.3\%$ a year. The chapter's 12.2\% also charges the hedges, whose sector
exposures change as positions open and close.
\end{solution}

\begin{solution}{exo:s1:residual-and-principal-component-stat-arb:5}
The residuals of a regression with an intercept sum to zero, so the cumulative residual starts and ends near zero. An AR(1) fitted to such a pinned path of
sixty points has a coefficient biased well below one, so nearly every path looks mean-reverting with a time far under thirty days, including a random walk.
The filter tests the estimate, and the estimate is biased.
\end{solution}

\begin{solution}{exo:s1:residual-and-principal-component-stat-arb:6}
The sector index removes the market and the stock's industry but not the styles; fifteen \omterm{def:s1:residual-and-principal-component-stat-arb:eigenportfolio}{eigenportfolios} span the market, the industries and the styles, so
the PCA residuals and the hedged positions carry less common risk: 7.5\% against 12.1\% of volatility with about the same number of positions.
\end{solution}

\begin{solution}{exo:s1:residual-and-principal-component-stat-arb:7}
1.22 at a share of 0.15 and $-0.06$ at zero. The second measures the method on a market with no mean-reverting level: its gross Sharpe ratio of 1.30 is the
planted one-day reversal leaking into the \omterm{def:s1:residual-and-principal-component-stat-arb:sscore}{s-score}, and costs of 9.8\% a year remove it. It is the method's floor.
\end{solution}

\begin{solution}{exo:s1:residual-and-principal-component-stat-arb:8}
The Sharpe ratio of 1.44 is an average over 1997--2007 that was much stronger before 2003 and 0.9 afterwards, and the cutoffs were chosen on part of the same
sample. The published parameters are not the risk; whether today's residuals revert is. Run it first on recent data, with the firm's costs, and measure the
mean-reverting share of residual variance directly.
\end{solution}

\begin{solution}{pb:s1:residual-and-principal-component-stat-arb:1}
\begin{enumerate}
\item Time-series regressions of each stock's last 60 returns on its sector fund (``ETF''), or on the returns of leading \omterm{def:s1:residual-and-principal-component-stat-arb:eigenportfolio}{eigenportfolios} (``PCA'').
\item $Q^{(j)}_i = v^{(j)}_i/\sigma_i$ from the correlation matrix's eigenvectors; 15, or as many as explain 55\% of the variance.
\item $\kappa = -252 \ln b$, $m = a/(1-b)$, $\sigma_{\mathrm{eq}} = \sqrt{\operatorname{Var}\zeta/(1-b^2)}$ from $X_{n+1} = a + bX_n + \zeta$.
\item $s = (X - m)/\sigma_{\mathrm{eq}} = -m/\sigma_{\mathrm{eq}}$; open at $\pm 1.25$, close longs at $-0.50$ and shorts at $0.75$.
\item A mean-reverting level in each stock's specific return taking a share of its variance, mean-reversion times around twenty days; the \omterm{def:s1:residual-and-principal-component-stat-arb:sscore}{s-score} correlates
$-0.128$ with the level's expected pull at a share of 0.3.
\item Sector: 0.07, 0.93, 2.15. PCA: $-0.06$, 1.22, 3.00.
\item The \omterm{def:s1:residual-and-principal-component-stat-arb:eigenportfolio}{eigenportfolios} remove the style factors as well; the PCA book's volatility is 7.5\% against 12.1\%.
\item Nothing after costs (0.07 and $-0.06$); its gross Sharpe ratio of 1.0--1.3 comes from the planted one-day reversal.
\item 99.3\% with, 99.1\% without: the AR(1) on a pinned 60-point path is biased toward fast reversion.
\item They reject more names and lower the net Sharpe ratio: 2.03, 1.81 and 0.74 for times under 15, 10 and 5 days.
\item Hardly: 2.20 to 2.30 between thresholds of 0.75 and 1.5, 1.87 at 2.0.
\item It lowers the net Sharpe ratio to 1.78: synthetic volume rises with the day's specific shock, so trading time shrinks the moves that reverse.
\item \textbf{Named result.} With a mean-reverting share of 0.3, the sector book nets 2.15 with the paper's speed filter and 2.16 without (PCA: 3.00 and
3.04); between entry thresholds of 0.75 and 1.5 the net ratio stays within 2.20--2.30, and at 2.0 it is 1.87.
\item The first eigenvector is nearly fixed (cosine 0.999 between monthly estimates); the fifteen-dimensional subspace overlaps 0.89 on average, and its least
stable direction only 0.63.
\item Sharpe ratios of 0.9 for PCA in 2003--2007, against 1.44 over 1997--2007; a similar degradation for the ETF version after 2002.
\item Books of this kind were unwound together, consistent with Khandani and Lo's account.
\item Residuals with a mean-reverting component large enough to pay the costs of trading a third of the book a day.
\item Measure the autocorrelation of residual returns at horizons of days to weeks on data the fits did not use, and the out-of-sample correlation of \omterm{def:s1:residual-and-principal-component-stat-arb:sscore}{s-scores}
with the next weeks' residual returns.
\item The cost model, and then the estimation window: the entry threshold and the speed filter hardly matter.
\item Trading stocks back toward where their factors say they should be, when they have strayed far enough to pay for the round trip.
\end{enumerate}
\end{solution}

\begin{solution}{iq:s1:residual-and-principal-component-stat-arb:1}
The distance of a stock's cumulative residual from its mean-reversion level, in equilibrium standard deviations, from an OU fit on a trailing window. Sell
(and buy the hedge) when it is high, buy when it is low, close when it comes back toward zero; size equally or through an optimiser.
\end{solution}

\begin{solution}{iq:s1:residual-and-principal-component-stat-arb:2}
PCA fits the market's actual common structure, including styles, and needs no fund data, but its factors rotate and are hard to interpret or hedge cheaply.
ETFs are tradable, stable and cheap to hedge, but leave style exposures in the residuals and depend on fund composition.
\end{solution}

\begin{solution}{iq:s1:residual-and-principal-component-stat-arb:3}
Not much. The AR(1) slope is biased downward in small samples, more so on a path pinned by the regression's intercept, and its standard error is large: a
mean-reversion time of 6 days ($b = 0.85$) is compatible with much slower reversion or none. Check with a bias-corrected estimate, a longer window, or the
fit's out-of-sample record.
\end{solution}

\begin{solution}{iq:s1:residual-and-principal-component-stat-arb:4}
Factor exposures left in the book (styles not in the hedge), a crowded unwind (other books with the same positions, as in August 2007), news in a few large
names, and the hedges: did the factors used for hedging change? Then decide whether to cut risk or provide liquidity.
\end{solution}

\begin{solution}{iq:s1:residual-and-principal-component-stat-arb:5}
Vectorise: one least-squares solve with all stocks as right-hand sides (the same factor matrix), cumulative sums and the AR(1) moments as column operations;
update the correlation matrix and eigenvectors weekly or monthly, not daily. That is a few matrix products per day.
\end{solution}

\begin{solution}{iq:s1:residual-and-principal-component-stat-arb:6}
The AR(1) $X_{n+1} = a + bX_n + \zeta$ has stationary variance $\operatorname{Var}\zeta/(1 - b^2)$, the square of $\sigma_{\mathrm{eq}}$. With an intercept in
the regression the residuals sum to zero, so $X_{60} = \sum_{n=1}^{60} \epsilon_n = 0$ and $s = (0 - m)/\sigma_{\mathrm{eq}}$.
\end{solution}
